\documentclass[12pt]{article}
\usepackage{amsmath, amsthm, amsfonts, mathtools, xcolor, caption}
\usepackage[margin=1in]{geometry}
\newtheorem{theorem}{Theorem}[section]
\newtheorem{lemma}[theorem]{Lemma}
\newtheorem{corollary}[theorem]{Corollary}
\newtheorem{proposition}[theorem]{Proposition}
\theoremstyle{definition}
\newtheorem{remark}{Remark}
\newtheorem{definition}{Definition}
\newtheorem{notation}{Notation}
\newtheorem{example}{Example}
\newtheorem{conjecture}[theorem]{Conjecture}
\newtheorem{openproblem}[theorem]{Open Problem}
\newcommand{\ev}[1]{( #1 )} 
\newcommand{\ew}[1]{\left( #1 \right)} 
\newcommand{\N}{\mathbb{N}}
\newcommand{\Z}{\mathbb{Z}}
\newcommand{\C}{\mathbb{C}}
\newcommand{\Q}{\mathbb{Q}}
\newcommand{\R}{\mathbb{R}}
\let\underbrace\LaTeXunderbrace
\let\overbrace\LaTeXoverbrace
\begin{document}
\section*{Challenge 6: Better-Quasi-Ordering of Finite Graphs under Minor Relation}
\refstepcounter{section}

\begin{definition}[Cumulative Hierarchy over $Q$]
Given a set $Q$, we denote by $\mathcal P^*$ the set of nonempty subsets of $Q$. We define for each ordinal $\alpha$ a set $V^*_{\alpha}(Q)$ recursively as follows:
\begin{itemize}
\item $V^*_0(Q) := Q$.
\item $V^*_{\beta + 1} := \cal{P}^* V^*_{\beta}$
\item $V^*_{\lambda}(Q) := \dot \bigcup_{\alpha < \lambda} V^*_{\alpha}(Q)$ for any limit ordinal $\lambda$.
\end{itemize}

The class given by the union of all these sets is denoted $V^*(Q)$. 
\end{definition}

\begin{definition}[$\alpha$-well-quasiordering]
Given a quasi-order $(Q,\le)$, we define a relation $\le^*$ on $V^*(Q)$ in terms of a 2-player game between Player I and Player II. A position in this game is a pair $(X,Y)$ with $X$ and $Y$ in $V_*(Q)$. If $X$ and $Y$ are both in $Q$ then this position is a win for Player II if $X \le Y$ and a win for Player I otherwise. If not, then Player I names a set $X'$, which must be $X$ itself if $X \in Q$ and otherwise must be an element of $X$. Then Player II also names a set $Y'$, which must be $Y$ itself if $Y \in Q$ and otherwise must be an element of $Y$. After this, the new position is $(X', Y')$. We define $X \le^* Y$ to mean that Player II has a winning strategy in this game with starting position $(X,Y)$.

We say that $(Q, \le)$ is {\em $\alpha$-well-quasiordered} if the restriction of this game to $V^*_{\alpha}(Q)$ is well-quasiordered.
\end{definition}

Note that $(Q, \le)$ is better-quasi-ordered if and only if it is $\omega_1$-well-quasi-ordered.

\begin{conjecture}
Let $r$ be an ordinal number. If $r$ is even, of the form $\alpha . 2$, then finite planar graphs are $\alpha$-well-quasi-ordered with respect to the minor relation. If $r$ is odd, of the form $\alpha . 2 + 1$, then finite graphs are $\alpha$-well-quasi-ordered.
\end{conjecture}

\end{document}